\chapter{Control Layer and Hardware Abstraction}
\label{ch:control-hal}

\section{Purpose of This Chapter}
\label{sec:control-hal-purpose}

This chapter documents how the \textbf{MICSBOL/ESP32\_BT\_Controller} firmware separates:

\begin{itemize}
    \item \textbf{control logic} (what the robot should do when the user moves a stick), and
    \item \textbf{hardware driving} (how the ESP32 writes PWM/GPIO/I\textsuperscript{2}C signals).
\end{itemize}

The repository implements this separation primarily through:

\begin{itemize}
    \item \texttt{control.cpp} / \texttt{control.h} --- the control layer
    \item \texttt{hal\_output.cpp} --- the hardware abstraction layer (HAL)
    \item \texttt{pins.h} --- the hardware mapping used by the HAL
    \item \texttt{pulse.cpp} / \texttt{pulse.h} --- non-blocking pulse outputs used by event actions
\end{itemize}

\section{Layered Architecture (High-Level View)}
\label{sec:control-hal-architecture}

The firmware uses a layered flow from the Android app to physical hardware:

\begin{enumerate}
    \item \textbf{Android app} transmits RC state packets and event packets over Bluetooth (SPP).
    \item \textbf{Receiver layer} (\texttt{receiver.cpp}) decodes packets and updates:
    \begin{itemize}
        \item \texttt{rcStatePacket} for continuous control values
        \item \texttt{rcEventPacket} for one-shot button events
    \end{itemize}
    \item \textbf{Control layer} (\texttt{control.cpp}) interprets \texttt{rcStatePacket} and calls HAL functions.
    \item \textbf{HAL} (\texttt{hal\_output.cpp}) performs the actual GPIO/PWM/I\textsuperscript{2}C operations.
\end{enumerate}


\begin{figure}[H]
    \centering
    \begin{tikzpicture}[
        font=\small,
        >=Stealth,
        line/.style={->, very thick, draw=black!70},
        box/.style={
            draw,
            rounded corners=2mm,
            very thick,
            align=center,
            minimum width=12.6cm,
            minimum height=1.15cm,
            fill=gray!4
        },
        boxApp/.style={box, fill=blue!6,  draw=blue!55!black},
        boxCore/.style={box, fill=teal!6, draw=teal!55!black},
        boxHw/.style={box, fill=green!6, draw=green!55!black},
        tag/.style={
            draw=black!25,
            fill=white,
            rounded corners=1mm,
            inner sep=2pt,
            font=\scriptsize\bfseries,
            text=black!70
        }
        ]

        % Manual coordinates (no "below=of" required)
        \node[boxApp]  (android) at (0, 0.0)  {Android App\\\footnotesize UI Controls (sticks, knobs, switches, buttons)};
        \node[boxCore] (bt)      at (0,-1.7)  {Bluetooth Serial (SPP)};
        \node[boxCore] (rx)      at (0,-3.4)  {receiver.cpp\\\footnotesize Rolling buffer + decode into rcStatePacket / rcEventPacket};
        \node[boxCore] (control) at (0,-5.1)  {control.cpp\\\footnotesize Control policy (map inputs to actions)};
        \node[boxCore] (hal)     at (0,-6.8)  {hal\_output.cpp\\\footnotesize Hardware outputs (PWM / GPIO / MCP)};
        \node[boxHw]   (hw)      at (0,-8.5)  {Hardware\\\footnotesize Motors, servos, LEDs, switches, MCP23017};

        % arrows
        \draw[line] (android) -- (bt);
        \draw[line] (bt) -- (rx);
        \draw[line] (rx) -- (control);
        \draw[line] (control) -- (hal);
        \draw[line] (hal) -- (hw);

        % left-side step tags
        \node[tag] at (-7.1, 0.0)  {1};
        \node[tag] at (-7.1,-1.7)  {2};
        \node[tag] at (-7.1,-3.4)  {3};
        \node[tag] at (-7.1,-5.1)  {4};
        \node[tag] at (-7.1,-6.8)  {5};
        \node[tag] at (-7.1,-8.5)  {6};

    \end{tikzpicture}
    \caption{Control data flow from Android to ESP32 hardware.}
\end{figure}

\section{Control Packets and State Structures}
\label{sec:control-hal-packets}

Control values are stored in \texttt{rcStatePacket} (defined in \texttt{packets.h}).
The fields include:

\begin{itemize}
    \item \texttt{leftStickX}, \texttt{leftStickY}
    \item \texttt{rightStickX}, \texttt{rightStickY}
    \item \texttt{leftKnob}, \texttt{rightKnob}
    \item \texttt{switches} (bit mask)
\end{itemize}

In the current firmware logic, stick and knob channels are treated as 12-bit values (nominal range 0 to 4095).

Button presses are delivered through \texttt{rcEventPacket} (also in \texttt{packets.h}) and are handled by
\texttt{controlHandleEvent(eventId)}.

\section{Control Layer (\texttt{control.cpp})}
\label{sec:control-hal-control-layer}

\subsection{Responsibilities}
\label{subsec:control-hal-control-resp}

The control layer is responsible for:

\begin{itemize}
    \item mapping sticks and knobs to actuator commands,
    \item selecting which HAL function to call for each logical output,
    \item decoding the switch bit mask into per-output states,
    \item handling discrete button events by triggering pulse outputs.
\end{itemize}

\subsection{Main loop: \texttt{controlUpdate()}}
\label{subsec:control-hal-controlUpdate}

The main control function is:

\begin{verbatim}
    controlUpdate()
\end{verbatim}

In the repository implementation, \texttt{controlUpdate()} reads values from \texttt{rcStatePacket} and calls:

\begin{itemize}
    \item \texttt{halSetSteering(...)}
    \item \texttt{halSetMotorLeft(...)}
    \item \texttt{halSetMotorRight(...)}
    \item \texttt{halSetCameraPan(...)}
    \item \texttt{halSetLed(...)}
    \item \texttt{halSetBuzzer(...)}
    \item \texttt{halSetSwitch(index, state)}
\end{itemize}

A key helper is the servo mapping function in \texttt{control.cpp}:

\begin{verbatim}
    static uint32_t mapServo(uint16_t v) {
        return map(v, 0, 4095, 205, 410);
    }
\end{verbatim}

\subsection{Switch bit mask decoding}
\label{subsec:control-hal-switch-mask}

Switches are transmitted as a single byte \texttt{rcStatePacket.data.switches}.
The control layer converts bits into outputs:

\begin{verbatim}
    for (int i = 0; i < 6; i++) {
        halSetSwitch(i, (sw >> i) & 1);
    }
\end{verbatim}

\subsection{Event handling: \texttt{controlHandleEvent(eventId)}}
\label{subsec:control-hal-events}

Discrete events (button presses) are handled by:

\begin{verbatim}
    controlHandleEvent(byte eventId)
\end{verbatim}

In the repository, events map to pulse outputs via \texttt{halPulseSwitch(index)}. The mapping depends on
project mode (direct GPIO vs MCP expander mode).

\section{Hardware Abstraction Layer (\texttt{hal\_output.cpp})}
\label{sec:control-hal-hal-layer}

\subsection{Responsibilities}
\label{subsec:control-hal-hal-resp}

The HAL implements the \emph{mechanisms} to drive hardware:

\begin{itemize}
    \item PWM output via ESP32 LEDC using \texttt{ledcWrite(...)}
    \item GPIO output via \texttt{digitalWrite(...)}
    \item MCP23017 expander output via \texttt{mcpWriteOutput(...)} (I\textsuperscript{2}C)
    \item motor direction pins and a dead-band mapping
    \item ramping motor speed using a periodic timer callback
    \item non-blocking pulse outputs using the pulse module
\end{itemize}

\subsection{PWM outputs (examples)}
\label{subsec:control-hal-hal-pwm}

Examples from \texttt{hal\_output.cpp}:

\begin{verbatim}
    void halSetSteering(uint32_t value) { ledcWrite(PIN_PWM_STICK_LX, value); }
    void halSetLed(uint32_t value)      { ledcWrite(PIN_PWM_KNOB_L, value);  }
    void halSetBuzzer(uint32_t value)   { ledcWrite(PIN_PWM_KNOB_R, value);  }
\end{verbatim}

\subsection{Motor control: direction + dead-band + ramping}
\label{subsec:control-hal-hal-motors}

Motor control in the HAL uses:

\begin{itemize}
    \item direction pins \texttt{PIN\_MD\_LEFT} and \texttt{PIN\_MD\_RIGHT}
    \item PWM output pins \texttt{PIN\_PWM\_STICK\_LY} (left) and \texttt{PIN\_PWM\_STICK\_RY} (right)
\end{itemize}

The HAL interprets motor values with a dead-band around the center, and it ramps actual motor PWM
toward a target value using a timer-driven update function.

\subsection{Switch outputs: direct GPIO vs MCP23017}
\label{subsec:control-hal-hal-switches}

The HAL provides a uniform interface:

\begin{verbatim}
    halSetSwitch(uint8_t index, bool state)
\end{verbatim}

but the implementation depends on build mode:

\begin{itemize}
    \item In direct GPIO mode, a small array of ESP32 pins is written using \texttt{digitalWrite(...)}.
    \item In MCP mode, the output is written using \texttt{mcpWriteOutput(...)} over I\textsuperscript{2}C.
\end{itemize}

\subsection{Pulse outputs (momentary actions)}
\label{subsec:control-hal-hal-pulses}

Pulse outputs are triggered by \texttt{halPulseSwitch(index)} and implemented using:

\begin{itemize}
    \item \texttt{pulseStart(pin, durationMs)} to begin a pulse
    \item \texttt{pulseUpdate()} to end it after the requested duration
\end{itemize}

This design is non-blocking, so the control loop can continue running while pulses are active.

\section{Pin Mapping (\texttt{pins.h}) and Portability}
\label{sec:control-hal-pins}

Hardware mappings are centralized in \texttt{pins.h}. When adapting the firmware to new hardware:

\begin{enumerate}
    \item update \texttt{pins.h} to match your wiring,
    \item adjust \texttt{hal\_output.cpp} if your motor driver or output logic differs,
    \item keep \texttt{control.cpp} stable as ``policy'' whenever possible.
\end{enumerate}

\section{Summary}
\label{sec:control-hal-summary}

The repository implements a clear separation between:

\begin{itemize}
    \item \textbf{Policy:} \texttt{control.cpp} interprets \texttt{rcStatePacket} and decides what actions to take.
    \item \textbf{Mechanism:} \texttt{hal\_output.cpp} drives PWM/GPIO/I\textsuperscript{2}C hardware to perform those actions.
\end{itemize}

This makes the firmware easier to port across different robots and wiring layouts, while keeping the codebase
organized and maintainable.
